% GENERATED FILE. Edit canonical lessons, then run npm run generate:books.
% canonical-source-hash: fnv1a64:f4a0503ef4daae9c
% canonical-lessons: TA-C542-uttaravatam, TA-C542-cellupati, TA-C542-muttirai, TA-C542-cattiyam-cey, TA-C542-ali-grant

\chapter{Guarantees and Oaths}
\label{ch:ta-guarantees-and-oaths}
\hlchaptermodality{542}

\begin{chapteropening}
\textbf{By the end of this chapter:} \emph{I can talk about guarantees, promises and stamps.}

Complete the last lesson of chapter 542: I can talk about guarantees, promises and stamps.
\end{chapteropening}

\section[\texorpdfstring{\ta{உத்தரவாதம்} (uttaravātam)}{உத்தரவாதம் (uttaravātam)}]{\ta{உத்தரவாதம்} (uttaravātam) — a guarantee; a warranty}

\label{lesson:TA-C542-uttaravatam}



\begin{quote}
Before the new one: say the Tamil for an agent, then the Tamil for bail.
\end{quote}

\subsection*{You'll want to know: \ta{உத்தரவாதம்}}
\textbf{\ta{உத்தரவாதம்}} — \emph{uttaravātam} — \textquotedblleft{}a guarantee; a warranty\textquotedblright{}.

\begin{quote}
\ta{இந்தக்} \ta{கைப்பேசிக்கு} \ta{ஒரு} \ta{வருட} \ta{உத்தரவாதம்} \ta{இருக்கிறது}.
\end{quote}

\begin{quote}
\emph{intak kaippēcikku oru varuṭa uttaravātam irukkiṟatu.}
\end{quote}

\begin{quote}
This phone has a one-year guarantee.
\end{quote}

With \ta{கொடு} (\emph{koṭu}): \ta{உத்தரவாதம்} \ta{கொடு}, to give a guarantee.

\begin{cousinweb}
\textbf{\ta{உத்தரவாதம்}} (\emph{uttaravātam}) is usually traced to Sanskrit \emph{uttara}, \textquotedblleft{}an answer, a reply\textquotedblright{}, and \emph{vāda}, \textquotedblleft{}a statement\textquotedblright{} (Monier-Williams): an answering for something.
\end{cousinweb}

\begin{grammarlens}[title={What you've built}]
One more word for saying what happens if.
\end{grammarlens}

\subsection*{Your turn}
\begin{itemize}
\raggedright
  \item \emph{Say it:} \emph{uttaravātam}
  \item \emph{Say it:} \emph{uttaravātam}, once more
  \item \emph{Say it:} \emph{intak kaippēcikku oru varuṭa uttaravātam irukkiṟatu.} — This phone has a one-year guarantee.
  \item \emph{Recall:} say the Tamil for an agent, then the Tamil for bail, then say \emph{uttaravātam} again
\end{itemize}

\begin{tcolorbox}[breakable,colback=teal!4,colframe=teal!35!black,title={Before you move on}]
What does \ta{உத்தரவாதம்} mean? (\textquotedblleft{}A guarantee; a warranty\textquotedblright{}.) Say it once more.
\end{tcolorbox}
\section[\texorpdfstring{\ta{செல்லுபடி} (cellupaṭi)}{செல்லுபடி (cellupaṭi)}]{\ta{செல்லுபடி} (cellupaṭi) — validity}

\label{lesson:TA-C542-cellupati}



\begin{quote}
Before the new one: say the Tamil for bail, then the Tamil for a guarantee.
\end{quote}

\subsection*{You'll want to know: \ta{செல்லுபடி}}
\textbf{\ta{செல்லுபடி}} — \emph{cellupaṭi} — \textquotedblleft{}validity\textquotedblright{}.

\begin{quote}
\ta{உங்கள்} \ta{சீட்டு} \ta{இன்று} \ta{வரை} \ta{மட்டுமே} \ta{செல்லுபடி} \ta{ஆகும்}.
\end{quote}

\begin{quote}
\emph{uṅkaḷ cīṭṭu iṉṟu varai maṭṭumē cellupaṭi ākum.}
\end{quote}

\begin{quote}
Your ticket is valid only until today.
\end{quote}

With \ta{ஆகு} (\emph{āku}) it says something is valid: \ta{செல்லுபடி} \ta{ஆகும்}.

\begin{cousinweb}
\textbf{\ta{செல்லுபடி}} (\emph{cellupaṭi}) is a native compound: \ta{செல்லு}, from \ta{செல்} (\emph{cel}), \textquotedblleft{}to go, to pass as coin, to be acceptable\textquotedblright{} (DEDR 2781), and \ta{படி}, \textquotedblleft{}a measure, a way\textquotedblright{}.
\end{cousinweb}

\begin{grammarlens}[title={What you've built}]
One more word for saying what happens if.
\end{grammarlens}

\subsection*{Your turn}
\begin{itemize}
\raggedright
  \item \emph{Say it:} \emph{cellupaṭi}
  \item \emph{Say it:} \emph{cellupaṭi}, once more
  \item \emph{Say it:} \emph{uṅkaḷ cīṭṭu iṉṟu varai maṭṭumē cellupaṭi ākum.} — Your ticket is valid only until today.
  \item \emph{Recall:} say the Tamil for bail, then the Tamil for a guarantee, then say \emph{cellupaṭi} again
\end{itemize}

\begin{tcolorbox}[breakable,colback=teal!4,colframe=teal!35!black,title={Before you move on}]
What does \ta{செல்லுபடி} mean? (\textquotedblleft{}Validity\textquotedblright{}.) Say it once more.
\end{tcolorbox}
\section[\texorpdfstring{\ta{முத்திரை} (muttirai)}{முத்திரை (muttirai)}]{\ta{முத்திரை} (muttirai) — a seal; an official stamp}

\label{lesson:TA-C542-muttirai}



\begin{quote}
Before the new one: say the Tamil for a guarantee, then the Tamil for validity.
\end{quote}

\subsection*{You'll want to know: \ta{முத்திரை}}
\textbf{\ta{முத்திரை}} — \emph{muttirai} — \textquotedblleft{}a seal; an official stamp\textquotedblright{}.

\begin{quote}
\ta{அலுவலர்} \ta{படிவத்தில்} \ta{முத்திரை} \ta{குத்தினார்}.
\end{quote}

\begin{quote}
\emph{aluvalar paṭivattil muttirai kuttiṉār.}
\end{quote}

\begin{quote}
The official stamped the form with a seal.
\end{quote}

One stamps with \ta{குத்து}: \ta{முத்திரை} \ta{குத்து}. A postage stamp is \ta{அஞ்சல்} \ta{தலை} (\emph{añcal talai}).

\begin{cousinweb}
\textbf{\ta{முத்திரை}} (\emph{muttirai}) is borrowed from Sanskrit \emph{mudrā}, \textquotedblleft{}a seal, a signet ring\textquotedblright{} (CDIAL 10203); Tamil wrote the \emph{dr} as \ta{த்திர}.
\end{cousinweb}

\begin{grammarlens}[title={What you've built}]
One more word for saying what happens if.
\end{grammarlens}

\subsection*{Your turn}
\begin{itemize}
\raggedright
  \item \emph{Say it:} \emph{muttirai}
  \item \emph{Say it:} \emph{muttirai}, once more
  \item \emph{Say it:} \emph{aluvalar paṭivattil muttirai kuttiṉār.} — The official stamped the form with a seal.
  \item \emph{Recall:} say the Tamil for a guarantee, then the Tamil for validity, then say \emph{muttirai} again
\end{itemize}

\begin{tcolorbox}[breakable,colback=teal!4,colframe=teal!35!black,title={Before you move on}]
What does \ta{முத்திரை} mean? (\textquotedblleft{}A seal; an official stamp\textquotedblright{}.) Say it once more.
\end{tcolorbox}
\section[\texorpdfstring{\ta{சத்தியம்} \ta{செய்} (cattiyam cey)}{சத்தியம் செய் (cattiyam cey)}]{\ta{சத்தியம்} \ta{செய்} (cattiyam cey) — to swear an oath}

\label{lesson:TA-C542-cattiyam-cey}



\begin{quote}
Before the new one: say the Tamil for validity, then the Tamil for a seal.
\end{quote}

\subsection*{You'll want to know: \ta{சத்தியம்} \ta{செய்}}
\textbf{\ta{சத்தியம்} \ta{செய்}} — \emph{cattiyam cey} — \textquotedblleft{}to swear an oath\textquotedblright{}.

\begin{quote}
\ta{சாட்சி} \ta{உண்மையைச்} \ta{சொல்வதாகச்} \ta{சத்தியம்} \ta{செய்தார்}.
\end{quote}

\begin{quote}
\emph{cāṭci uṇmaiyaic colvatākac cattiyam ceytār.}
\end{quote}

\begin{quote}
The witness swore to tell the truth.
\end{quote}

Past (I …): \ta{சத்தியம்} \ta{செய்தேன்}, \emph{cattiyam ceytēṉ}.

What one swears is said with -\ta{ஆக} after a verbal noun: \ta{வருவதாகச்} \ta{சத்தியம்} \ta{செய்}, to swear to come.

\begin{cousinweb}
\textbf{\ta{சத்தியம்} \ta{செய்}} (\emph{cattiyam cey}) joins \ta{சத்தியம்} (\emph{cattiyam}), \textquotedblleft{}truth, an oath\textquotedblright{}, borrowed from Sanskrit \emph{satya}, \textquotedblleft{}true; a promise\textquotedblright{} (CDIAL 13112), and the light verb \ta{செய்} (\emph{cey}), \textquotedblleft{}to do\textquotedblright{}.
\end{cousinweb}

\begin{grammarlens}[title={What you've built}]
One more word for saying what happens if.
\end{grammarlens}

\subsection*{Your turn}
\begin{itemize}
\raggedright
  \item \emph{Say it:} \emph{cattiyam cey}
  \item \emph{Say it:} \emph{cattiyam cey}, once more
  \item \emph{Say it:} \emph{cāṭci uṇmaiyaic colvatākac cattiyam ceytār.} — The witness swore to tell the truth.
  \item \emph{Recall:} say the Tamil for validity, then the Tamil for a seal, then say \emph{cattiyam cey} again
\end{itemize}

\begin{tcolorbox}[breakable,colback=teal!4,colframe=teal!35!black,title={Before you move on}]
What does \ta{சத்தியம்} \ta{செய்} mean? (\textquotedblleft{}To swear an oath\textquotedblright{}.) Say it once more.
\end{tcolorbox}
\section[\texorpdfstring{\ta{அளி} (aḷi)}{அளி (aḷi)}]{\ta{அளி} (aḷi) — to grant; to give formally}

\label{lesson:TA-C542-ali-grant}



\begin{quote}
Before the new one: say the Tamil for a seal, then the Tamil for to swear an oath.
\end{quote}

\subsection*{You'll want to know: \ta{அளி}}
\textbf{\ta{அளி}} — \emph{aḷi} — \textquotedblleft{}to grant; to give formally\textquotedblright{}.

\begin{quote}
\ta{அரசு} \ta{ஏழைக்} \ta{குடும்பங்களுக்கு} \ta{உதவி} \ta{அளித்தது}.
\end{quote}

\begin{quote}
\emph{aracu ēḻaik kuṭumpaṅkaḷukku utavi aḷittatu.}
\end{quote}

\begin{quote}
The government gave help to poor families.
\end{quote}

Past (I …): \ta{அளித்தேன்}, \emph{aḷittēṉ}.

A formal verb beside \ta{கொடு} (\emph{koṭu}). The receiver takes -\ta{க்கு}.

\begin{cousinweb}
\textbf{\ta{அளி}} (\emph{aḷi}) is native Dravidian: DEDR 301 gives \ta{அளி} (\emph{aḷi}), \textquotedblleft{}to protect, to show favour\textquotedblright{}, and a Malayalam cognate that means \textquotedblleft{}to grant, to bestow\textquotedblright{}, the sense the Tamil verb has today.
\end{cousinweb}

\begin{grammarlens}[title={What you've built}]
One more word for saying what happens if.
\end{grammarlens}

\subsection*{Your turn}
\begin{itemize}
\raggedright
  \item \emph{Say it:} \emph{aḷi}
  \item \emph{Say it:} \emph{aḷi}, once more
  \item \emph{Say it:} \emph{aracu ēḻaik kuṭumpaṅkaḷukku utavi aḷittatu.} — The government gave help to poor families.
  \item \emph{Recall:} say the Tamil for a seal, then the Tamil for to swear an oath, then say \emph{aḷi} again
\end{itemize}

\begin{tcolorbox}[breakable,colback=teal!4,colframe=teal!35!black,title={Before you move on}]
What does \ta{அளி} mean? (\textquotedblleft{}To grant; to give formally\textquotedblright{}.) Say it once more.
\end{tcolorbox}
